\documentclass[12pt,a4paper]{article}
\usepackage[english]{babel}
\usepackage{titlesec}

% \titleformat{\chapter} % command [display] % shape {\bfseries\huge} % format
%   {Chapter \thechapter} % label {0pt} % sep {} % before-code


\usepackage[T5]{fontenc}
\usepackage{mathptmx}[ptm]
\usepackage{a4wide, amssymb, epsfig, latexsym, array, hhline, fancyhdr}
\usepackage[normalem]{ulem}
%\usepackage{soul}
\usepackage{float}
\usepackage[makeroom]{cancel}
\usepackage{amsmath}
\usepackage{amsthm}
\usepackage{multicol, longtable, amscd}
\usepackage{diagbox}%Make diagonal lines in tables
\usepackage{booktabs}
\usepackage{alltt}
\usepackage[framemethod=tikz]{mdframed}% For highlighting paragraph backgrounds
\usepackage{indentfirst}
\usepackage{caption,subcaption}
\usepackage{svg}
\graphicspath{ {figures/} }
\usepackage{lastpage}
\usepackage[lined,boxed,commentsnumbered]{algorithm2e}
\usepackage{enumerate}
\usepackage{color}
\usepackage{graphicx}% Standard graphics package
\graphicspath{ {figures/} }
\usepackage{array}
\usepackage{tabularx, caption}
\usepackage{tabularray}
\usepackage{multirow}
\usepackage{multicol}
\usepackage{rotating}
\usepackage{graphics}
\usepackage{geometry}
\usepackage{setspace}
\usepackage{epsfig}
\usepackage{tikz}
\usepackage[numbers]{natbib}
\usepackage{subcaption}
\usepackage{dirtree}

\usetikzlibrary{arrows, snakes, backgrounds, calc}
\usepackage[unicode]{hyperref}
\hypersetup{urlcolor=blue,linkcolor=black,citecolor=black,colorlinks=true} 

\usepackage{xcolor}

\definecolor{codegreen}{rgb}{0,0.6,0}
\definecolor{codegray}{rgb}{0.5,0.5,0.5}
\definecolor{codepurple}{rgb}{0.58,0,0.82}
\definecolor{backcolour}{rgb}{0.95,0.95,0.92}

\usepackage{listings}
\lstdefinestyle{mystyle}{ backgroundcolor=\color{backcolour},   
    commentstyle=\color{codegreen}, keywordstyle=\color{magenta},
    numberstyle=\tiny\color{codegray}, stringstyle=\color{codepurple},
    basicstyle=\ttfamily\footnotesize, breakatwhitespace=false,         
    breaklines=true,                 
    captionpos=b,                    
    keepspaces=true,                 
    numbers=left,                    
    numbersep=5pt,                  
    showspaces=false,                
    showstringspaces=false, showtabs=false,                  
    tabsize=2 }
\lstset{style=mystyle}



\def\thesislayout{	% A4: 210 × 297
	\geometry{ a4paper, total={160mm,240mm},  % fix over page
		left=30mm, top=30mm, } }
\thesislayout

%\usepackage{fancyhdr}
\setlength{\headheight}{40pt}
\pagestyle{fancy}
\fancyhead{} % clear all header fields
\fancyhead[L]{
 \begin{tabular}{rl}
    \begin{picture}(25,15)(0,0) \put(0,-8){\includegraphics[width=10mm,
    height=10mm]{Images/hcmut.png}}
    %\put(0,-8){\epsfig{width=10mm,figure=hcmut.eps}}
   \end{picture}&
	%\includegraphics[width=8mm, height=8mm]{hcmut.png} & %
	\begin{tabular}{l}
        \textbf{ Ho Chi Minh City University of Technology } \\
        \textbf{ Faculty of Computer Science and Engineering }
	\end{tabular} 	
 \end{tabular}
} \fancyhead[R]{
	\begin{tabular}{l}
		\tiny \bf \\
		\tiny \bf \end{tabular}  } \fancyfoot{} % clear all footer fields
\fancyfoot[L]{\scriptsize Modern Speech Processing (CO5257) - Academic year 2024 - 2025}
\fancyfoot[R]{\scriptsize  Page {\thepage}/\pageref{LastPage}}
\renewcommand{\headrulewidth}{0.3pt}
\renewcommand{\footrulewidth}{0.3pt}

%%%%%%%%%%%%%%%%%%%%%%%%%%%%%%%%%%%%%%%%%%
% Tạo một môi trường mới cho các bảng đặc tả UC, giúp định dạng kiểu
\newenvironment{usecase_table}[1][Test]
{
\begin{table}[H]
% Left-Right Table padding
    \setlength{\tabcolsep}{6pt}
    \renewcommand{\arraystretch}{1.5}
    % \begin{tabularx}{\textwidth}{bs}
    \def\savedcaption{\caption{#1}}%
    \begin{tabular}{|>{\bfseries} p{0.2\linewidth}|p{0.7\linewidth}|}

} {
    \end{tabular}
    \savedcaption
\end{table}
}

\newenvironment{usecase_enum}
{
    \begin{enumerate*}[itemjoin={\newline}]
} {
    \end{enumerate*}
}

\makeatletter
\newenvironment{acknowledgments}{
	\small
	\begin{center}
	  {\bfseries Acknowledgments\vspace{-.5em}\vspace{\z@}}
	\end{center}
	\quotation
  }{}
\makeatother

% Declaration
\makeatletter
\newenvironment{declaration}{
	\small
	\begin{center}
	  {\bfseries Declaration\vspace{-.5em}\vspace{\z@}}
	\end{center}
	\quotation
  }{}
\makeatother

\makeatletter
\newenvironment{abstr}{
        \small
        \begin{center}
	{\bfseries Abstract \vspace{-.5em}\vspace{\z@}}
        \end{center}
        \quotation
    }{}
\makeatother

\makeatletter
\newenvironment{abstrsss}{
        \small
        \begin{center}
	{\bfseries Abstractss \vspace{-.5em}\vspace{\z@}}
        \end{center}
        \quotation
    }{}
\makeatother
%%%%%%%%%%%%%%%%%%%%%%%%%%%%%%%%%%%%%%%%%%%%%%5
\setcounter{secnumdepth}{4}
\setcounter{tocdepth}{3}
\makeatletter
\newcounter {subsubsubsection}[subsubsection]
\renewcommand\thesubsubsubsection{\thesubsubsection .\@alph\c@subsubsubsection}
\newcommand\subsubsubsection{\@startsection{subsubsubsection}{4}{\z@}%
                                     {-3.25ex\@plus -1ex \@minus -.2ex}%
                                     {1.5ex \@plus .2ex}%
                                     {\normalfont\normalsize\bfseries}}
\newcommand*\l@subsubsubsection{\@dottedtocline{3}{10.0em}{4.1em}}
\newcommand*{\subsubsubsectionmark}[1]{}
\makeatother



\everymath{\color{black}}%make in-line maths symbols blue to read/check easily

\sloppy
\captionsetup[figure]{labelfont={small,bf},textfont={small,it},belowskip=-1pt,aboveskip=-9pt}
%space removed between caption, figure, and text
\captionsetup[table]{labelfont={small,bf},textfont={small,it},belowskip=-1pt,aboveskip=7pt}
\setlength{\floatsep}{5pt plus 2pt minus 2pt}
\setlength{\textfloatsep}{5pt plus 2pt minus 2pt}
\setlength{\intextsep}{10pt plus 2pt minus 2pt}


% Declare a variable named Proc using \newcommand

\newcommand{\Uni}{Ho Chi Minh University - Vietnam National Unversity, Ho Chi Minh City}

\renewcommand{\labelitemi}{\normalfont\bfseries\textendash}
\renewcommand{\labelitemii}{\textbullet}
\renewcommand*\baselinestretch{1.5}\selectfont


\thesislayout

\begin{document}

\begin{titlepage}
    \begin{tikzpicture}[remember picture, overlay]
        \draw[line width = 4pt] ($(current page.north west) + (1.0in,-0.5in)$)
        rectangle ($(current page.south east) + (-0.4in,0.5in)$);
        \draw[line width=1.5pt]
        ($ (current page.north west) + (1.05in,-0.55in) $) rectangle ($ (current
        page.south east) + (-0.45in,0.55in) $);
    \end{tikzpicture}

    \begin{center}
        \large \textbf{VIETNAM NATIONAL UNIVERSITY, HO CHI MINH CITY} \\
        \large \textbf{HO CHI MINH UNIVERSITY OF TECHNOLOGY} \\
        \large \textbf{FACULTY OF COMPUTER SCIENCE AND ENGINEERING}
    \end{center}

    \begin{figure}[h!]
        \begin{center}
            \includegraphics[width=5cm]{Images/hcmut.png}
        \end{center}
    \end{figure}

    \begin{center}
        \begin{tabular}{c}
        \multicolumn{1}{l}{\textbf{{\Large Modern Speech Processing (CO5257)}}}\\
        ~~\\
        \hline
        \\
        \multicolumn{1}{l}{\textbf{{\Large Assignment Report}}}\\
        \\
        \textbf{{\Large Utilizing Generative AI for Context-Aware Note
        Generation}}\\
        \textbf{{\Large in Real-Time Speech Processing}}\\
        \\
        \hline
        \end{tabular}
        \end{center}

    \vspace{1cm}
    \begin{table}[H]
        \begin{tabular}{rrl}
        \hspace{5 cm} & \textbf{Mentor}: & Nguyễn Đức Dũng\\
        
        & \textbf{Student}: & Trần Hà Tuấn Kiệt -- 2011493 \\
        
        \end{tabular}
        \end{table}
    \vspace{1cm}

    \begin{center}
        {\large Ho Chi Minh City, \the\month/\the\year}
    \end{center}
\end{titlepage}

\newpage

%\thispagestyle{empty} \chapter*{Work Assignment} \begin{table}[H] \centering
% \renewcommand{\arraystretch}{1.5} \begin{tabular}{|c|c|c|l|} \hline
% \textbf{No.}       & \textbf{Full name}                 & \textbf{Student ID}
% & \textbf{Work Assignment}                     \\
%         \hline %%%%%Student 1%%%%%%%%%%% \multirow{4}{*}{1} &
%         \multirow{4}{*}{Trần Hà Tuấn Kiệt} & \multirow{4}{*}{2011493} & -
%         Exploratory Data Analysis                  \\
%                            &                                    &
%                            & - System proposal and architecture           \\
%                            &                                    &
%                            & - Research and Develop with new technologies \\
%                            &                                    &
%                            & - Implement system endpoints                 \\
%         \hline %%%%%Student 2%%%%%%%%%%% \multirow{4}{*}{2} &
%         \multirow{4}{*}{Nguyễn Đức Thuỵ}   & \multirow{4}{*}{2012158} & -
%         Exploratory Data Analysis                  \\
%                            &                                    &
%                            & - System proposal and architecture           \\
%                            &                                    &
%                            & - Research and Develop with new technologies \\
%                            &                                    &
%                            & - Implement data processing flows            \\
%         \hline \end{tabular} \end{table} \newpage

% Revert title name back to English
\renewcommand{\contentsname}{Table of Contents}
\renewcommand{\listfigurename}{List of Figures}
\renewcommand{\listtablename}{List of Tables}
\renewcommand{\bibname}{References}
\renewcommand{\figurename}{Figure}

\tableofcontents
\newpage

%\thispagestyle{empty} \renewcommand\thechapter{Chapter \arabic{chapter}}

\section*{Abstract}
In the era of information overload, efficient and accurate note-taking during
real-time speech interactions is paramount across various sectors, including
education, business, and healthcare. Traditional method often fall short in
capturing and filtering all context of discussions, leading to incomplete,
unrelated or misinterpreted notes. This research explores the potential of
Generative Artificial Intelligence (AI) to develop a context-aware note taking
system capable of processing and summarizing real-time speech inputs. By
integrating natural language processing (NLP) and automatic speech recognition
(ASR) to transcribe interactions, coupled with advanced prompting techniques to
generate draft notes using large language models (LLMs), the system dynamically
interprets the semantic nuances of ongoing conversations, ensuring that the
generated notes accurately reflect the intended meaning and key points
discussed. Preliminary evaluations demonstrate the system's proficiency in
maintaining contextual integrity and relevance in the notes produced, thereby
enhancing comprehension and retention for users. The findings suggest
significant potential for deploying such AI-driven solutions to augment human
capabilities in environments where real-time information processing is critical.

\newpage
\section{Introduction}
In today's information-rich environment, effective and precise note-taking
during real-time speech interactions is crucial across various sectors,
including education, business, and healthcare. Accurate documentation at the
point of care enhances patient safety by ensuring timely and precise recording
of clinical information, thereby reducing errors in treatment. %WIKIPEDIA 
Moreover, the adoption of electronic health records and point-of-care
documentation devices has streamlined workflows, allowing healthcare providers
to capture data directly during patient encounters, which minimizes redundant
tasks and improves communication among care teams. % WIKIPEDIA 
Similarly, in educational settings, structured note-taking methods like the
Cornell Notes system have been shown to aid in the retention and comprehension
of lecture material, thereby enhancing learning outcomes\cite{boyle2014note}.
These examples underscore the importance of efficient note-taking practices in
managing information effectively across different professional domains.

Traditional methods of note-taking have been further enhanced by technological
advancements, particularly through automatic speech recognition (ASR) systems.
These systems rely on two primary components: an acoustic model (AM), which
deciphers the relationship between audio signals and phonetic units, and a
language model (LM), which predicts the likelihood of word sequences to ensure
grammatical coherence. The language model is typically trained independently
using extensive corpora of transcribed speech collected by existing ASR systems.
Despite their effectiveness, traditional ASR systems often face challenges in
capturing the full context of interactions \cite{itsp2022}.

Generative AI, a subset of artificial intelligence that creates new content
based on existing data, builds on these foundations to offer significant
opportunities for enhancing context-aware documentation practices. By
integrating advanced natural language processing (NLP) techniques with ASR
capabilities, generative AI systems can produce coherent and contextually
relevant documentation, thereby improving efficiency and accuracy across various
sectors. This convergence of ASR and generative AI heralds a new era in
real-time documentation, addressing limitations of traditional systems and
paving the way for more dynamic and adaptive solutions.

This research focuses on the development and evaluation of generative AI systems
for context-aware note generation during real-time speech interactions. The
study aims to bridge the gap between traditional ASR systems and advanced
generative models by incorporating contextual understanding to produce coherent,
accurate, and relevant notes. By leveraging techniques such as attention-based
mechanisms, multi-modal data integration, and real-time natural language
processing, the research seeks to address challenges like omissions,
inaccuracies, and the inability to adapt to complex scenarios. The scope
includes testing these models in diverse applications such as clinical
documentation, educational environments, and business meetings, highlighting
their potential to transform how notes are generated and utilized across various
domains.

\newpage
\section{Previous Work}
The integration of generative AI into real-time speech processing has led to
significant advancements in context-aware note generation. Notable contributions
include the development of frameworks that combine large language models (LLMs),
retrieval-augmented generation (RAG), and automatic speech recognition (ASR) to
automate the creation of structured notes from medical conversations. These
systems capture and process both text and voice inputs, generating contextually
accurate documentation \cite{leong2024gen}.

% Whisper
The development of Whisper (Alec .et al. 2022\cite{radford2022robust})
represents a significant advancement in speech recognition technology. Trained
on an extensive dataset of 680,000 hours of multilingual and multitask
supervised data, Whisper demonstrates robust generalization across various
benchmarks, often matching the performance of fully supervised models in
zero-shot transfer scenarios without requiring fine-tuning. Notably, its
accuracy and resilience are comparable to human capabilities, marking a
substantial step forward in the field. The release of Whisper's models and
inference code provides a solid foundation for future research and development
in robust speech processing.

Kernberg et al. (2023) \cite{biswas2024intelligent} conducted an evaluation of
ChatGPT-4, a conversational AI developed by OpenAI utilizing GPT-3.5 and GPT-4
large language models, to assess its effectiveness in generating SOAP
(Subjective, Objective, Assessment, and Plan) notes from transcripts of
simulated patient-provider interactions. The study identified an average of 23.6
errors per clinical case, with omissions constituting 86\% of these errors.
Additionally, there was notable variability between different iterations of the
same case, with only 52.9\% of data elements consistently reported across all
three replicates. The accuracy of the AI-generated notes was inversely
correlated with both the length of the transcripts and the complexity of the
data elements, indicating challenges in managing complex medical scenarios. The
authors concluded that the quality and reliability of the AI-generated clinical
notes did not meet the standards necessary for clinical application, emphasizing
the need for further research to address issues related to accuracy,
variability, and potential errors.

In the realm of speech recognition, integrating contextual information has
proven to enhance the performance of neural language models (NLMs). Martinez et
al. (2021) \cite{martinez2021attention} introduced an attention-based mechanism
that incorporates both textual and non-linguistic contextual data into NLMs,
resulting in a 7.0\% relative reduction in perplexity over standard language
models lacking contextual information. Similarly, Chang et al. (2021)
\cite{chang2021context} developed a Context-Aware Transformer Transducer (CATT)
that leverages multi-head attention to integrate contextual signals, achieving a
24.2\% improvement in word error rate compared to baseline transformer
transducer models. These advancements underscore the efficacy of attention-based
methods in enhancing the adaptability and accuracy of speech recognition systems
through the incorporation of contextual cues.

Furthermore, hybrid systems that combine generative AI with domain-specific
ontologies and knowledge graphs have demonstrated significant potential in
improving note generation. These systems leverage structured data to provide
contextual augmentation to language models, resulting in outputs that are both
more accurate and domain-relevant. For instance, Wang et al. (2023)
\cite{wang2023ontology} explored the use of knowledge graph-enhanced
transformers to generate real-time summaries in legal and academic domains,
achieving a notable reduction in redundancy and an increase in the relevance of
generated notes.

In medical settings, the integration of ASR with electronic health records (EHR)
systems has facilitated the seamless transfer of generated notes into clinical
workflows. Projects like Ambient Clinical Intelligence (ACI) by companies such
as Nuance and 3M have demonstrated the feasibility of embedding generative AI
into practice, enabling physicians to focus more on patient care while the AI
handles documentation \cite{nuance2023aci}. Despite these advancements,
challenges such as ensuring data privacy, handling domain-specific jargon, and
achieving a balance between automation and human oversight remain critical areas
of ongoing research.

The application of generative AI in education is another emerging area, with
systems designed to transcribe lectures, summarize discussions, and generate
personalized study notes for students. Tools like Otter.ai and Scribe
\cite{otter2023scribe} utilize a combination of ASR and generative summarization
techniques to create concise and coherent notes from complex discussions.
Studies by Chen et al. (2023) \cite{chen2023education} highlight the utility of
such systems in improving students' comprehension and retention, especially in
multidisciplinary and high-information-density courses.

Lastly, ethical considerations play a pivotal role in the deployment of
generative AI for note generation. Issues such as bias in language models, the
potential for misinformation, and user consent in voice data processing require
careful scrutiny. Researchers are actively exploring frameworks to ensure
fairness, transparency, and accountability in generative AI systems. Recent
efforts by Binns et al. (2024) \cite{binns2024ethical} propose a set of
guidelines for designing equitable and trustworthy AI systems, which include
stakeholder engagement, periodic model audits, and the adoption of interpretable
AI techniques.

Collectively, these advancements and challenges illustrate the multifaceted
nature of generative AI in context-aware note generation, highlighting both its
transformative potential and the need for continued innovation and ethical
vigilance. 
\newpage

\section {Methodology}

The methodology employed in this research revolves around the integration of
generative AI with ASR and natural language processing (NLP) techniques to
generate contextually relevant notes during real-time speech interactions. The
study followed a systematic workflow, encompassing data collection and
processing, natural language understanding, contextual awareness, prompt
engineering, model selection and deployment.



\section {Data Collection and Processing}

\newpage
\section {Natural Language Understanding}

\newpage
\section {Contextual Awareness}

\newpage
\section {Prompt Engineering}

\newpage
\section {Model Selection and Deployment}

\newpage
\section{Experimental Setup}

\section{Evaluations}

\newpage
\section{Limitations and Further Research}

\newpage
\nocite{*}
\bibliographystyle{plainnat}

\bibliography{refs}

\end{document}